%%%%%%%%%%%%%%%%%%%%%%%%%%%%%%%%%%%%%%%%%%%%%%%%%%%%%%%%%%%%%%%%%%%%%%%%%%%%%%%
%%  ** NÃO ALTERAR ESTE ARQUIVO **
%%  Processamento dos dados da pasta *info/*: aqui ficam os valores derivados
%%  (nomes completos, título conferido, descrição etc.). Os alunos editam
%%  apenas os arquivos de info/. Carregado por principal.tex, ata.tex e ficha.tex,
%%  sempre logo depois de %%%%%%%%%%%%%%%%%%%%%%%%%%%%%%%%%%%%%%%%%%%%%%%%%%%%%%%%%%%%%%%%%%%%%%%%%%%%%%%
%%  ** NÃO ALTERAR ESTE ARQUIVO **
%%  Reúne os dados do trabalho, que o aluno edita nos arquivos desta pasta:
%%    trabalho.tex   fase, tipo, curso e ano
%%    orientacao.tex orientador, coorientador e departamento
%%    autores.tex    autores (alunos)
%%    projeto.tex    título, subtítulo e palavras-chave
%%    banca.tex      membros da banca
%%    ficha.tex      dados da ficha catalográfica
%%%%%%%%%%%%%%%%%%%%%%%%%%%%%%%%%%%%%%%%%%%%%%%%%%%%%%%%%%%%%%%%%%%%%%%%%%%%%%%
\newcommand{\faseAtual}{Projeto} % [Projeto|Final]
\newcommand{\tipo}{Conclusão de Curso} % [Iniciação Científica|Conclusão de Curso|Mestrado|Doutorado]
\newcommand{\Curso}{Ciência da Computação} % Aterar para o seu curso se necessário
\newcommand{\ano}{2026} % Alterar para o ano atual

% Dados do Orientador
% Titulação: Dr | Me | (vazio). Gênero: vazio para orientador, "a" para orientadora.
\newcommand{\orientadorNome}{Fagner de Assis}
\newcommand{\orientadorSobrenome}{Moura Pimentel}
\newcommand{\orientadorTitulo}{Dr}
\newcommand{\orientadorGenero}{}
% Coorientador (opcional: deixe o nome vazio se não houver).
\newcommand{\coorientadorNome}{}
\newcommand{\coorientadorSobrenome}{}
\newcommand{\coorientadorTitulo}{}
\newcommand{\coorientadorGenero}{}
\newcommand{\departamento}{Ciência da Computação}

% Autores (alunos): um \autor{Nome}{Sobrenome} por autor, na ordem da capa e da ficha.
% O primeiro é o autor principal da ficha catalográfica.
\newcommand{\autores}{
    \autor{Nome}{Sobrenome do Autor 1} % Alterar
    \autor{Nome}{Sobrenome do Autor 2} % Alterar, adicionar ou remover
}

% Dados do Projeto
\newcommand{\titulo}{Título do Trabalho \todo{alterar}}
\newcommand{\sub}{Subtítulo do Trabalho \todo{alterar}}

% Palavras-chave (de 3 a 5 termos, separados por vírgula).
% Usadas no resumo (pretextuais/resumo.tex) e na ficha catalográfica (ficha.tex).
\newcommand{\listaPalavrasChave}{Palavra-chave um, Palavra-chave dois, Palavra-chave três, Palavra-chave quatro} % Alterar

% Banca (além do orientador): um \membroBanca{Nome} por membro, na ordem das assinaturas da ata.
\newcommand{\banca}{
    \membroBanca{Prof. Dr. Fulano} % Alterar
    \membroBanca{Prof. Dr. Beltrano} % Alterar, adicionar ou remover
}

% Ficha catalográfica (não é usado em IC)
% Os campos seguem o formulário da biblioteca: https://ficha.fei.edu.br/ficha/
% Título e subtítulo vêm de \titulo e \sub (info/projeto.tex); deixe \sub vazio se não houver.
\newcommand{\fichaCidade}{São Bernardo do Campo}
\newcommand{\fichaContagem}{f} % [f = folhas|p = páginas]
\newcommand{\fichaNumero}{20} % Alterar: número de folhas ou páginas
\newcommand{\fichaIlustracoes}{sim} % [sim|vazio]: possui ilustrações
\newcommand{\fichaFonte}{Times} % [Times|Arial]
\newcommand{\fichaTamanho}{10} % [9|10]

.
%%%%%%%%%%%%%%%%%%%%%%%%%%%%%%%%%%%%%%%%%%%%%%%%%%%%%%%%%%%%%%%%%%%%%%%%%%%%%%%
\RequirePackage{etoolbox}

% --- Autores: \alunos ("Nome Sobrenome \\ Nome Sobrenome") a partir de \autores ---
% Também guarda cada autor (\autorKn = nome, \autorKs = sobrenome) e o total (\numAutores),
% que a ficha catalográfica usa.
\newcount\numAutores \numAutores=0
\def\alunos{}
\newcommand{\autor}[2]{\advance\numAutores by 1
  \expandafter\gdef\csname autor\the\numAutores n\endcsname{#1}%
  \expandafter\gdef\csname autor\the\numAutores s\endcsname{#2}%
  \ifnum\numAutores>1 \appto\alunos{\\}\fi
  \appto\alunos{#1 #2}}
\autores
\appto\alunos{\\} % termina com quebra de linha, como esperado pela capa e pela ata

% --- Banca: \assinaturasBanca (uma linha de assinatura por membro) a partir de \banca ---
\def\assinaturasBanca{}
\newcommand{\membroBanca}[1]{\appto\assinaturasBanca{\hrulefill \\ #1\par\vspace*{0.05\textheight}}}
\banca

% --- Orientador: "Prof.ª Dra. ", "Prof. Me. " e nome completo ---
\newcommand{\tratamentoProf}[2]{Prof.\ifdefempty{#2}{}{\textordfeminine}\ %
  \ifdefstring{#1}{Dr}{\ifdefempty{#2}{Dr.}{Dra.}\ }{%
  \ifdefstring{#1}{Me}{\ifdefempty{#2}{Me.}{Ma.}\ }{}}}
\newcommand{\orientador}{\tratamentoProf{\orientadorTitulo}{\orientadorGenero}\orientadorNome\ \orientadorSobrenome}

% --- Derivados de \tipo: título conferido (\diploma) e tipo na ficha (\fichaTipo) ---
\ifdefstring{\tipo}{Mestrado}
  {\newcommand{\diploma}{Mestre}\newcommand{\fichaTipo}{Dissertação}}
  {\ifdefstring{\tipo}{Doutorado}
    {\newcommand{\diploma}{Doutor}\newcommand{\fichaTipo}{Tese}}
    {\ifdefstring{\tipo}{Iniciação Científica}
      {\newcommand{\diploma}{}\newcommand{\fichaTipo}{Iniciação Científica}}
      {\newcommand{\diploma}{Bacharel}\newcommand{\fichaTipo}{Trabalho de Conclusão de Curso}}}}

% --- Texto da folha de rosto e da ata ---
% Iniciação Científica não tem a parte "como parte dos requisitos... título de ... em ...".
\ifdefstring{\tipo}{Iniciação Científica}
  {\newcommand{\descricao}{Trabalho de \tipo~ apresentado ao Centro Universitário FEI. Orientad\ifdefempty{\orientadorGenero}{o}{a} por \orientador.}}
  {\newcommand{\descricao}{Trabalho de \tipo~ apresentado ao Centro Universitário FEI, como parte dos requisitos necessários para obtenção do título de \diploma~ em \Curso . Orientad\ifdefempty{\orientadorGenero}{o}{a} por \orientador.}}

% --- Palavras-chave do resumo: "A; B; C" a partir de \listaPalavrasChave ---
\newcommand{\palavraItem}[1]{\palavraSep #1\gdef\palavraSep{; }}
\newcommand{\palavrasChaveResumo}{\gdef\palavraSep{}\expandafter\forcsvlist\expandafter\palavraItem\expandafter{\listaPalavrasChave}}
